% Part: second-order-logic
% Chapter: metatheory
% Section: compactness

\documentclass[../../../include/open-logic-section]{subfiles}

\begin{document}

\olfileid[hi]{sol}{met}{com}

\olsection{द्वितीय-कोटि तर्कशास्त्र संहत नहीं है}

\begin{explain}
!!{sentence} के किसी समुच्चय~$\Gamma$ को \emph{परिमिततः संतुष्ट्य} कहें यदि
उसका प्रत्येक परिमित उपसमुच्चय संतुष्ट्य हो। प्रथम-कोटि तर्कशास्त्र में यह
गुण है कि यदि !!{sentence} का कोई समुच्चय~$\Gamma$ परिमिततः संतुष्ट्य है,
तो वह संतुष्ट्य है। इस गुण को \emph{संहतता} कहते हैं। अनुगमन की भाषा में
इसका एक तुल्य रूप है: यदि $\Gamma \Entails !A$, तो किसी परिमित
उपसमुच्चय~$\Gamma_0 \subseteq \Gamma$ के लिए पहले ही $\Gamma_0 \Entails
!A$। इस रूप में यह पूर्णता प्रमेय का तात्कालिक उपप्रमेय है: यदि $\Gamma
\Entails !A$, तो पूर्णता से $\Gamma \Proves !A$। किंतु कोई !!{derivation},
~$\Gamma$ के केवल परिमित रूप से अनेक !!{sentence} का ही उपयोग कर सकता है।

द्वितीय-कोटि तर्कशास्त्र के लिए संहतता सत्य नहीं है। द्वितीय-कोटि
!!{sentence} के ऐसे समुच्चय हैं जो परिमिततः संतुष्ट्य तो हैं, किंतु संतुष्ट्य
नहीं; और जो किसी~$!A$ का अनुगमन करते हैं, जबकि उनका कोई परिमित उपसमुच्चय
उस~$!A$ का अनुगमन नहीं करता।
\end{explain}


\begin{thm}
\ollabel{thm:sol-undecidable}
द्वितीय-कोटि तर्कशास्त्र संहत नहीं है।
\end{thm}

\begin{proof}
स्मरण करें कि
\[
\fn{Inf} \ident \lexists[u][(\lforall[x][\lforall[y][(\eq[u(x)][u(y)]
      \lif \eq[x][y])]] \land \lexists[y][\lforall[x][\eq/[y][u(x)]]])]
\]
किसी !!a{structure} में तभी और केवल तभी संतुष्ट होता है जब उसका प्रांत
अनंत हो। $!A^{\ge n}$ ऐसा वाक्य हो जो कहता है कि प्रांत में कम-से-कम
$n$ !!{element} हैं; उदाहरणार्थ,
\[
!A^{\ge n} \ident \lexists[x_1][\dots\lexists[x_n][(\eq/[x_1][x_2]
    \land \eq/[x_1][x_3] \land \dots \land \eq/[x_{n-1}][x_n])]].
\]
!!{sentence} के इस समुच्चय पर विचार करें:
\[
\Gamma = \{\lnot \fn{Inf}, !A^{\ge 1}, !A^{\ge 2}, !A^{\ge 3}, \dots\}.
\]
यह परिमिततः संतुष्ट्य है, क्योंकि किसी भी परिमित उपसमुच्चय~$\Gamma_0
\subseteq \Gamma$ के लिए कोई ऐसा $k$ है कि $!A^{\ge k} \in \Gamma$,
पर $n>k$ के लिए कोई $!A^{\ge n} \in \Gamma$ नहीं है। यदि $\Domain{M}$
में $k$ !!{element} हों, तो $\Sat{M}{\Gamma_0}$। किंतु $\Gamma$ संतुष्ट्य
नहीं है: यदि $\Sat{M}{\lnot \fn{Inf}}$, तो $\Domain{M}$ परिमित होना चाहिए,
मान लें उसका आकार~$k$ है। तब $\Sat/{M}{!A^{\ge k+1}}$।
\end{proof}

\begin{prob}
ऐसे समुच्चय~$\Gamma$ और !!a{sentence}~$!A$ का उदाहरण दें कि
$\Gamma \Entails !A$, पर प्रत्येक परिमित उपसमुच्चय~$\Gamma_0 \subseteq
\Gamma$ के लिए $\Gamma_0 \Entails/ !A$।
\end{prob}


\end{document}
